\makeabstract
    {
    Depolymerization of Waste PMMA
    }
    {
    Eduardo Ramirez-Ventura\textsuperscript{1}, Tanmoy Maity\textsuperscript{2}, Jacob M. Mirsky\textsuperscript{2}, Bryan Cowley\textsuperscript{2}, Brent. S. Sumerlin\textsuperscript{2}
    }
    {
    \textsuperscript{1} Department of Chemistry, Amherst College, Amherst, MA\\
\textsuperscript{2} Department of Chemistry, University of Florida, Gainesville, FL
    }
    {
    Poly(methyl methacrylate) (PMMA) is widely used for its mechanical properties, but high thermodynamic and kinetic barriers make depolymerization difficult, leading to inefficient waste recycling. Recent work has been done on pre-made polymers but this study aims to modify pre-existing PMMA waste. By incorporating phthalimide ester pendant groups our goal is to achieve high percentages of PMMA depolymerization at lower temperatures.
    }
    {
    Our approach begins with partial hydrolysis (\ensuremath{<}15\%) of colored PMMA waste to install pendent carboxyl groups. Afterwards esterification with N-hydroxyphthalimide to create phthalimide-ester containing PMMA. The phthalimide-ester groups provide thermally sensitive triggers that allow for effective depolymerization at lower temperatures. Nuclear Magnetic Resonance (1H NMR) spectra was used to confirm conversion of pendent groups and Thermogravimetric Analysis (TGA) was used to determine extent of depolymerization.
    }
    {
    1H NMR spectra suggests hydrolysis of colored PMMA and subsequent conversion to N-hydroxyphthalimide ester. TGA ramp suggested that depolymerization onset decreased to 240{\textdegree}C significantly lower than PMMA. Isothermal hold showed 85\% mass loss, equivalent to 85\% depolymerization. Successful post-polymerization modification of P(MMA-co-MAA) and high extents of depolymerization of PMMA were achieved.
    }
    {
    We showed that our approach is effective for depolymerizing real world PMMA waste and helps contribute to a circular polymer economy. Future work will focus on expanding depolymerization to other vinyl polymers and post-automotive waste. More work should also be done to scale depolymerization toward techno-economic viability. 
    }

    \newpage
    